% =====================================================================
%  CHAPTER 11: 健康危险因素评价（对应大纲第十一章）
%  参考PPT：第十一章 健康危险因素评价
% =====================================================================
\chapter{健康危险因素评价}
\setcounter{chapter}{11}
\chapterintro{
掌握健康危险因素的概念、分类及特点；掌握健康危险因素评价的概念和步骤（难点*）；掌握危险分数的意义和来源（难点*）；掌握评价年龄和增长年龄的概念；掌握健康危险因素的个体评价和群体评价方法。
}

\section{健康危险因素概述}

\begin{defbox}
\textbf{健康危险因素（Health Risk Factors）：}是指在机体内外环境中存在的与疾病发生、发展和死亡有关的诱发因素，即能使疾病或死亡发生的可能性增加的因素。

\textbf{分类：}环境因素（自然环境+社会环境）、生物遗传因素、行为生活方式因素、卫生服务因素。

\textbf{特点：}潜伏期长、特异性弱、联合作用明显、广泛存在。

\textbf{健康危险因素的作用过程：}危险因素→机体→疾病→健康结局。
\end{defbox}

\section{健康危险因素评价}

\begin{keybox}
\textbf{健康危险因素评价（Health Risk Factors Appraisal, HRA）：}是研究危险因素与慢性疾病发病率及死亡率之间数量依存关系及其规律性的一种技术，研究人们生活在危险因素的环境中发生死亡或发病的概率，以及当改变不良行为、消除或降低危险因素时，死亡及发病危险降低的程度。

\textbf{评价步骤（难点*）：}
\begin{enumerate}[leftmargin=*]
    \item \imp{收集资料}：(A)当地性别、年龄别疾病死亡率资料（作为平均危险水平）；(B)个人健康危险因素资料（通过问卷/体检/实验室检查等收集）。一般选择10$\sim$15种主要死亡原因作为研究对象
    \item \imp{将危险因素转换成危险分数}：危险分数是衡量某项危险因素对健康危害程度的指标。基准值为1.0——=1.0表示与人群平均水平相同，>1.0表示危险度高于平均水平，<1.0表示危险度低于平均水平
    \item \imp{计算组合危险分数}：多个危险因素综合作用时的危险分数（注意协同作用而非简单相加）
    \item \imp{计算评价年龄}：评价年龄 = 根据个人存在的各项危险因素计算出的总死亡危险相当于一般人群中多大年龄者的死亡危险水平
    \item \imp{计算增长年龄}：增长年龄 = 将个人存在的可改变危险因素去除后计算出的总死亡危险相当于的年龄
    \item \imp{计算危险降低程度}：危险降低程度 = (评价年龄 $-$ 增长年龄) / 评价年龄 × 100\%
\end{enumerate}

\textbf{个体评价方法：}
\begin{itemize}[leftmargin=*]
    \item 评价年龄>日历年龄：存在危险因素
    \item 评价年龄<日历年龄：危险因素低于平均水平
    \item 增长年龄<评价年龄：有降低危险的潜力
\end{itemize}

\textbf{群体评价方法：}根据人群危险因素分布特征和危险程度，确定高危人群，评价干预效果。
\end{keybox}

\section{复习题}

\begin{enumerate}[leftmargin=*, itemsep=8pt]
    \item \textbf{【名词解释】}健康危险因素；健康危险因素评价（HRA）；危险分数；评价年龄；增长年龄
    \item \textbf{【简答题】}简述健康危险因素的概念、分类及特点。
    \item \textbf{【简答题】}试述健康危险因素评价的六个步骤（难点*）。
    \item \textbf{【简答题】}简述危险分数的意义和来源。
    \item \textbf{【简答题】}简述健康危险因素的个体评价方法和群体评价方法。
\end{enumerate}

\subsection*{参考答案要点}

\begin{enumerate}[leftmargin=*, itemsep=5pt]
    \item \textbf{健康危险因素}：机体内外环境中存在的与疾病发生/发展和死亡有关的诱发因素。\textbf{HRA}：研究危险因素与慢性疾病发病率及死亡率之间数量依存关系及其规律性的技术。\textbf{危险分数}：衡量某项危险因素对健康危害程度的指标，基准值1.0，>1危险度高，<1危险度低。\textbf{评价年龄}：根据个人危险因素计算出的总死亡危险相当于一般人群中多大年龄者的水平。\textbf{增长年龄}：去除可改变危险因素后计算出的总死亡危险相当于的年龄。
    \item 四大类：环境因素（自然+社会）、生物遗传因素、行为生活方式因素、卫生服务因素。四大特点：潜伏期长、特异性弱、联合作用明显、广泛存在。
    \item (1)收集资料（A.当地性别年龄别疾病死亡率 B.个人危险因素资料，选10$\sim$15种主要死因）；(2)危险因素→危险分数；(3)计算组合危险分数；(4)计算评价年龄；(5)计算增长年龄；(6)计算危险降低程度。
    \item 危险分数是衡量某项危险因素对健康危害程度的指标，基准值为1.0（与人群平均水平相同）。>1.0危险度高于平均水平，<1.0低于平均水平。来源：通过流行病学研究（队列研究为主）获得的相对危险度转换而来。
    \item 个体评价：比较评价年龄与日历年龄、增长年龄与评价年龄。群体评价：根据人群危险因素分布特征和危险程度确定高危人群，评价干预效果。
\end{enumerate}

\newpage
